\documentclass[11pt]{article}

\usepackage[a4paper,margin=1in]{geometry}
\usepackage{amsmath,amssymb}
\usepackage{booktabs}
\usepackage{microtype}
\usepackage{hyperref}
\usepackage{xcolor}

\hypersetup{colorlinks=true,linkcolor=blue!55!black,urlcolor=blue!55!black}

\newcommand{\RR}{\mathbb R}
\newcommand{\VolZ}{\operatorname{Vol}_{\mathbb Z}}
\newcommand{\MV}{\operatorname{MV}}

\title{Supplementary Computational Certificate for the $8216$ Three-Zone Bound}
\author{Chaoyu Hu \and Tianrun Yu}
\date{}

\begin{document}
\maketitle

\begin{abstract}
This supplement records the exact finite computation used in the Laurent--BKK stage of the accompanying paper.  The committed JSON file contains every exponent vector, every Newton vertex, and the fifteen normalized Minkowski-sum volumes for each of three polynomial systems.  A deterministic verifier reconstructs the data from the Laurent formulas and proves
\[
 \MV(\mathcal P_0)=6808,
 \qquad
 \MV(\mathcal P_1)=1568,
 \qquad
 \MV(\mathcal P_2)=1248.
\]
These values give $6808+(1568+1248)/2=8216$.
\end{abstract}

\section{Files and verification command}

The certificate consists of
\begin{center}
\begin{tabular}{@{}p{0.56\textwidth}p{0.36\textwidth}@{}}
\toprule
File & Contents\\
\midrule
\nolinkurl{verification/verify_certificate.py}
  & exact symbolic, support, hull, and volume verifier\\
\nolinkurl{verification/mixed_volume_certificate.json}
  & stable machine-readable certificate\\
\bottomrule
\end{tabular}
\end{center}
It is checked by
\begin{verbatim}
python3 verification/verify_certificate.py \
  --check verification/mixed_volume_certificate.json
\end{verbatim}
The command regenerates the JSON and compares it byte for byte with the committed file.  It reports
\begin{verbatim}
support/vertex counts: F=36/16, G=265/24, H2=1154/52,
H1_tilde=12360/54, P_epsilon=13/8, Hc_tilde=2692/74
normalized polarization sums: 163392, 37632, 29952
mixed volumes: 6808, 1568, 1248
regular positive three-zone bound: 8216
\end{verbatim}

\section{Exact Laurent support construction}

All monomials have the form
\[
 z^aw^bR^cS^d,
 \qquad (a,b,c,d)\in\mathbb Z^4.
\]
The verifier implements sparse Laurent arithmetic over an exact formal coefficient ring generated by
\[
 \alpha,\ \beta,\ p,\ c_-,\ \overline{c_-},\
 c_+,\ \overline{c_+},\ \frac{\mu^2}{16B^2}.
\]
The physical involution sends $z,w,p$ to their inverses, exchanges each $c_\pm$ with its formal conjugate, and fixes the real parameters.  Thus all coefficient cancellations are decided exactly, without numerical parameter substitution.

After substituting the three displayed flight numerators from the paper, sparse division in $R$ verifies
\[
 (R-1)^3\mid F_0,
 \qquad
 (R-1)^4\nmid F_0,
 \qquad
 (R-1)^2\mid G_0.
\]
The quotients $F_L=F_0/(R-1)^3$ and $G_L=G_0/(R-1)^2$ generate the two basic supports.  Contact supports are then obtained solely by union, integer Minkowski addition, derivative translation, and the determinant formulas in the paper.  In particular, the covering-contact set is generated directly rather than replaced by a rectangular envelope.

\begin{center}
\begin{tabular}{@{}lrr@{}}
\toprule
Object & exponent vectors & Newton vertices\\
\midrule
$F_L$ & $36$ & $16$\\
$G_L$ & $265$ & $24$\\
$H_2$ & $1154$ & $52$\\
$\widetilde H_1$ & $12360$ & $54$\\
$P_\varepsilon$ & $13$ & $8$\\
$\widetilde H_c$ & $2692$ & $74$\\
\bottomrule
\end{tabular}
\end{center}
The JSON keys \texttt{supports} and \texttt{vertices} contain the corresponding complete lexicographically ordered lists.

\section{Exact hull and volume checks}

SciPy/Qhull is used only to propose hull vertices and a triangulation of each boundary.  For every proposed nondegenerate facet, the verifier computes a primitive integer normal from cofactors and checks that all vertices lie in one closed supporting half-space.  It chooses the rational barycenter of the vertices as an interior cone point.  Every cone-simplex determinant is recomputed using \texttt{Fraction}; their absolute values sum to
\[
 \VolZ(P)=4!\operatorname{Vol}_{\mathrm{Eucl}}(P).
\]
Degenerate zero-volume simplices introduced by a triangulation of a nonsimplicial rational facet are discarded exactly.  An accepted normalized volume is required to be an integer.

For four polytopes $P_1,\ldots,P_4$, the BKK convention in the paper is
\begin{equation}\label{eq:polarization}
 \MV(P_1,P_2,P_3,P_4)=\frac1{4!}
 \sum_{\varnothing\ne I\subseteq\{1,2,3,4\}}
 (-1)^{4-|I|}\VolZ\left(\sum_{i\in I}P_i\right).
\end{equation}
The factor $1/4!$ converts normalized lattice volumes back to the coefficient convention of Bernstein's theorem.

\section{The forty-five Minkowski-sum volumes}

In each row below the subsets occur in mask order
\[
 1,\ 2,\ 12,\ 3,\ 13,\ 23,\ 123,\
 4,\ 14,\ 24,\ 124,\ 34,\ 134,\ 234,\ 1234.
\]
The entries are normalized volumes $\VolZ(P_I)$.

\begin{center}
\small
\begin{tabular}{@{}lrrrrr@{}}
\toprule
System & $1$ & $2$ & $12$ & $3$ & $13$\\
\midrule
$\mathcal P_0$ & 192 & 2240 & 14784 & 14780 & 50364\\
$\mathcal P_1$ & 192 & 2240 & 14784 & 64 & 2528\\
$\mathcal P_2$ & 192 & 2240 & 14784 & 14780 & 50364\\
\bottomrule
\end{tabular}

\medskip
\begin{tabular}{@{}lrrrrr@{}}
\toprule
System & $23$ & $123$ & $4$ & $14$ & $24$\\
\midrule
$\mathcal P_0$ & 104508 & 236540 & 207530 & 414450 & 640426\\
$\mathcal P_1$ & 12480 & 46496 & 39760 & 107264 & 200752\\
$\mathcal P_2$ & 104508 & 236540 & 64 & 2528 & 12480\\
\bottomrule
\end{tabular}

\medskip
\begin{tabular}{@{}lrrrrr@{}}
\toprule
System & $124$ & $34$ & $134$ & $234$ & $1234$\\
\midrule
$\mathcal P_0$ & 1099938 & 1099878 & 1757342 & 2396278 & 3553822\\
$\mathcal P_1$ & 405504 & 95912 & 217912 & 379592 & 695672\\
$\mathcal P_2$ & 46496 & 46476 & 120620 & 225868 & 445612\\
\bottomrule
\end{tabular}
\end{center}

Substitution into \eqref{eq:polarization} gives
\begin{align*}
 \mathcal P_0:&\quad 163392/24=6808,\\
 \mathcal P_1:&\quad 37632/24=1568,\\
 \mathcal P_2:&\quad 29952/24=1248.
\end{align*}

\section{Certified implication}

The two-step Khovanskii--Rolle inequality in the paper is
\[
 N_{\mathrm{reg},3z}^{+}
 \le \MV(\mathcal P_0)
 +\frac12\bigl(\MV(\mathcal P_1)+\MV(\mathcal P_2)\bigr).
\]
The exact values above therefore prove the finite computational implication
\[
 \boxed{N_{\mathrm{reg},3z}^{+}\le
 6808+\frac12(1568+1248)=8216}.
\]

\end{document}
